\subsection{RQ3: Admission Decision Quality}
\label{sec:rq3}

Table~\ref{tab:rq3_admission_audit} compares our approach with naive approaches on the same audited decisions.
To assess decision quality independently of runtime enforcement, we audited 3{,}746 decisions on the S26 Ultra under an always-admit configuration.
Optional stages executed with the watchdog suppressed.
Admission decisions were recorded but not enforced.
Each decision is thus scored against its realized outcome as feasible if its remaining sequence completed within the remaining timeout budget at that decision, and unsafe otherwise.

\begin{table}[t]
  \centering
  \setlength{\abovecaptionskip}{2pt}%
  \setlength{\belowcaptionskip}{3pt}%
  \caption{RQ3: Admission decision quality against naive approaches.}
  \label{tab:rq3_admission_audit}

  {\scriptsize
  \renewcommand{\arraystretch}{1.12}%
  \newlength{\rqnfawidth}\settowidth{\rqnfawidth}{\textbf{97.3\% (3{,}470)}}%
  \newlength{\rqnfswidth}\settowidth{\rqnfswidth}{\textbf{4.4\% (157)}}%
  \newlength{\rqnuawidth}\settowidth{\rqnuawidth}{23.6\% (42)}%
  \newlength{\rqnuswidth}\settowidth{\rqnuswidth}{\textbf{97.2\% (173)}}%
  \newcommand{\rqnfa}[1]{\makebox[\rqnfawidth][r]{#1}}%
  \newcommand{\rqnfs}[1]{\makebox[\rqnfswidth][r]{#1}}%
  \newcommand{\rqnua}[1]{\makebox[\rqnuawidth][r]{#1}}%
  \newcommand{\rqnus}[1]{\makebox[\rqnuswidth][r]{#1}}%
  \newsavebox{\rqnbox}%

  \newcommand{\rqntabular}{%
  \begin{tabular}{@{}%
      >{\raggedright\arraybackslash}p{86pt}|%
      cc|cc%
    @{\hspace{\tabcolsep}}}
    \toprule
    \multirow[c]{3}{=}[-1.6pt]{\centering\textbf{Admission approach}}
      & \multicolumn{4}{c}{\textbf{Admission audit, \% (\(n\))}} \\
    \cmidrule(l){2-5}
      & \multicolumn{2}{c|}{\textbf{Feasible decisions}}
      & \multicolumn{2}{c}{\textbf{Unsafe decisions}} \\
    \cmidrule(lr){2-3}\cmidrule(l){4-5}
      & \multicolumn{1}{c}{\textbf{Admitted}}
      & \multicolumn{1}{c|}{\textbf{Skipped}}
      & \multicolumn{1}{c}{\textbf{Admitted}}
      & \multicolumn{1}{c}{\textbf{Skipped}} \\
    \midrule

    \multicolumn{5}{@{}l}{\textbf{Naive approaches}} \\
    \addlinespace[1pt]
    Mean of last 3 stage times & \rqnfa{99.6\% (3{,}554)} & \rqnfs{0.4\% (14)} & \rqnua{23.6\% (42)} & \rqnus{76.4\% (136)} \\
    Max of last 3 stage times  & \rqnfa{98.2\% (3{,}505)} & \rqnfs{1.8\% (63)} & \rqnua{7.9\% (14)} & \rqnus{92.1\% (164)} \\
    Max of all prior stage times & \rqnfa{93.7\% (3{,}344)} & \rqnfs{6.3\% (224)} & \rqnua{2.2\% (4)} & \rqnus{97.8\% (174)} \\
    \midrule

    \textbf{Our approach}
      & \rqnfa{\textbf{97.3\% (3{,}470)}} & \rqnfs{\textbf{2.7\% (98)}}
      & \cellcolor{gray!20}\rqnua{\textbf{2.8\% (5)}} & \rqnus{\textbf{97.2\% (173)}} \\
    \hspace{8pt}12MP normal, \(M\)       & \rqnfa{98.1\% (831)} & \rqnfs{1.9\% (16)} & \rqnua{\textbf{5.7\% (2)}} & \rqnus{94.3\% (33)} \\
    \hspace{8pt}12MP normal, \(S\)       & \rqnfa{99.4\% (842)} & \rqnfs{0.6\% (5)} & \rqnua{\textbf{0.0\% (0)}} & \rqnus{100.0\% (35)} \\
    \hspace{8pt}24MP memory, \(M\)       & \rqnfa{94.7\% (892)} & \rqnfs{5.3\% (50)} & \rqnua{\textbf{1.9\% (1)}} & \rqnus{98.1\% (53)} \\
    \hspace{8pt}24MP memory, \(S\)       & \rqnfa{97.1\% (905)} & \rqnfs{2.9\% (27)} & \rqnua{\textbf{3.7\% (2)}} & \rqnus{96.3\% (52)} \\
    \bottomrule
  \end{tabular}}%
  \fittabcolsep{\rqnbox}{\rqntabular}{\columnwidth}{10}%
  \rqntabular
  }
\end{table}

\input{figures/fig_rq3_unsafe_spike_anatomy.tex}

\myparagraph{Decision accuracy}
Among the evaluated approaches, only our approach kept both feasible-skip and unsafe-admit rates below 3\% in the pooled audit.
It skipped 98 of 3{,}568 feasible decisions (2.7\%) and admitted 5 of 178 unsafe decisions (2.8\%).
The naive approaches that sum the mean or maximum of each remaining stage's last three durations admitted 42 and 14 unsafe decisions, respectively.
The naive approach summing each remaining stage's maximum prior duration admitted four unsafe decisions but skipped 224 feasible decisions, 2.3 times as many as our approach.

\myparagraph{Unsafe admits}
The five unsafe admits coincided with latency increases not captured by the decision-time estimates.
Figure~\ref{fig:rq3_unsafe_spike_anatomy} compares the remaining-sequence latency and CPU time of these five decisions with those of the preceding capture.
M1--M3 and S1--S2 label the multi-frame and single-frame decisions, respectively.
For M1 and M2, latency rose to 2.32\(\times\) and 1.95\(\times\), whereas CPU time rose only to 1.41\(\times\) and 1.15\(\times\).
The average number of busy cores, calculated as CPU time divided by latency, consequently fell to 0.60\(\times\) and 0.59\(\times\).

\myparagraph{Safeguard replay}
Trace replay identified a safeguard response for all five unsafe admits using the recorded stage durations, budgets, upper estimates, and admission decisions.
For M1 and M2, stage durations exceeded watchdog budgets (1{,}460 vs.\ 1{,}202~ms and 1{,}648 vs.\ 889~ms, respectively), so the replayed watchdog triggered fallback before the stages returned.
For M3, the recorded decision for a later \(S\) stage in the same capture was a skip.
Without that stage, the replayed sequence fit within the budget.
For S1 and S2, the recorded \(M\) decisions earlier in the same captures were skips.
Without \(M\), the replayed budget at the \(S\) decision covered the remaining sequence.

